\section{Exercise 3}
\textit{Define the potential $\phi$ as:}
\[\phi(x,y,z) = ln(x^2 + y^2 + z^2), where (x,y,z) \in \mathbb{R}^3 \{(0,0,0)\}\]

\begin{itemize}
    \item Calculate $\nabla\phi$
    \item Calculate the flux of $\nabla\phi$ out of the sphere: $S^2 = \{(x,y,z)\in\mathbb{R}^3 |x^2+y^2+z^2 = R^2\}$ where R is a positive real number.
\end{itemize}

The gradient is defined as the partial derivative of each of the of the variables. In this case the potential $\phi$ is
is a logarithmic function, meaning it has to be treated as an inner and an outer function, using the chain rule.
$f(x) = ln(g(x)) \Rightarrow \frac{d}{dx} f(x) \Rightarrow f'(x) = \frac{1}{g(x)} \cdot g(x)$\\
\[\frac{\partial \phi}{\partial x} = 
\frac{\partial}{\partial x}ln(x^2 + y^2 + z^2) = 
\frac{1}{x^2+y^2+z^2} \cdot (2x + 0 + 0) = 
\frac{2x}{x^2+y^2+z^2}\]
\[\frac{\partial \phi}{\partial y} = 
\frac{\partial}{\partial y}ln(x^2 + y^2 + z^2) = 
\frac{1}{x^2+y^2+z^2} \cdot (0 + 2y + 0) = 
\frac{2y}{x^2+y^2+z^2}\]
\[\frac{\partial \phi}{\partial z} = 
\frac{\partial}{\partial z}ln(x^2 + y^2 + z^2) = 
\frac{1}{x^2+y^2+z^2} \cdot (0 + 0 + 2z) = 
\frac{2z}{x^2+y^2+z^2}\]
Yielding the final result:
\[\nabla\phi = \left(\frac{2x}{x^2+y^2+z^2}, \frac{2y}{x^2+y^2+z^2},\frac{2z}{x^2+y^2+z^2} \right)\]
To calculate the flux out of the unit sphere $S^2$, this yields another scalar field, which is $\nabla \cdot \nabla$, this
is also known as the \textit{laplacian} of $\phi$. This is possible if the scalar field is twice differentiable, which this is. 
And then afterwards integrate over the volume enclosed by the sphere.
\[\nabla\cdot\nabla\phi = \nabla^2 \phi = 
\frac{\partial }{\partial x}\left(\frac{\partial \phi}{\partial x}\right) + 
\frac{\partial }{\partial y}\left(\frac{\partial \phi}{\partial y}\right) +
\frac{\partial }{\partial z}\left(\frac{\partial \phi}{\partial z}\right)
\]
Deriving the laplacian:
\[\nabla^2\phi = 
\frac{\partial }{\partial x}\left(\frac{2x}{x^2+y^2+z^2}\right) +
\frac{\partial }{\partial y}\left(\frac{2y}{x^2+y^2+z^2}\right) +
\frac{\partial }{\partial z}\left(\frac{2z}{x^2+y^2+z^2}\right)
\]
This is evalutated using the quotiend rule:$(\frac{u}{v})' = \frac{u'v - uv'}{v^2}$
\[\frac{\partial}{\partial x}\left( \frac{2x}{x^2+y^2+z^2} \right) = \frac{2(x^2+y^2+z^2) - 4x^2}{(x^2+y^2+z^2)^2}\]
The same procedure is done to derive the partial derivation expression for y and x, finally yielding the expressiom
\[\nabla^2\phi\ = \frac{2(x^2+y^2+z^2) - 4x^2 - 2(x^2+y^2+z^2) -4y^2 + 2(x^2+y^2+z^2) - 4z^2}{(x^2+y^2+z^2)^2}\]
\[= \frac{6(x^2+y^2+z^2) - 4(x^2+y^2+z^)}{(x^2+y^2+z^2)^2}\]
Cleaning up the expression:
\[= \frac{2}{x^2+y^2+z^2}\]
And by the definition of the sphere, this can be expressed as:
\[\nabla\phi =  \frac{2(x^2+y^2+z^2)}{x^2+y^2+z^2} = \frac{2}{R^2}\]
Flux through the sphere can then be evalutated as:
\[\Phi = \iiint_{S^2} \frac{2}{R^2}dV\]
Using the same bounds of integration as before:
\[\Phi = \int_{0}^{1}\int_{0}^{\pi}\int_{0}^{2\pi} 2sin\theta \ d\phi \ d\theta \  dr\]
first integral:
\[2sin\theta \int_{0}^{2\pi}d\phi = 2sin\theta \left[\phi\right]^{2\pi}_0 = 2sin\theta(2\pi - 0) = 4\pi sin\theta\]
second integral:
\[\int_{0}^{\pi}4\pi sin\theta \ d\theta = 4\pi\int_{0}^{\pi} sin\theta \ d\theta = 4\pi \left[-cos\theta\right]^\pi_0 = 
4\pi \ ((-cos\pi) - (-cos(0))) = 4\pi \ (-(-1) - (-1)) = 4\pi \ (1 + 1) = 8\pi\]
third integral:
\[\int_{0}^{1}8\pi dr = 8\pi \int_{0}^{1}dr = 8\pi \ \left[r\right]^1_0 = 8\pi \ (1 - 0) = 8\pi \]
This results in the flux being:
\[\Phi = 8\pi\]